\documentclass[11pt,a4paper]{article}

% --- Encoding and language ---------------------------------------------
% fontenc/inputenc are safe defaults for pdfLaTeX; XeLaTeX ignores them.
\usepackage[T1]{fontenc}
\usepackage[utf8]{inputenc}

% --- Core packages ------------------------------------------------------
\usepackage{amsmath, amssymb}   % maths
\usepackage{graphicx}           % \includegraphics
\usepackage{booktabs}           % better table rules
\usepackage[hidelinks]{hyperref}% clickable refs without coloured boxes

\title{Your Paper Title}
\author{Your Name\thanks{Your Institution, \texttt{you@example.com}}}
\date{\today}

\begin{document}
\maketitle

\begin{abstract}
  One paragraph summarising the problem, what you did, and what you found.
  Write it last.
\end{abstract}

\section{Introduction}
\label{sec:introduction}

State the problem and why it matters. Cite prior work as you go, for example
the classic reference~\cite{knuth1984}.

\section{Method}
\label{sec:method}

Describe what you did in enough detail that a reader could repeat it.
Equations are numbered automatically and can be referenced:

\begin{equation}
  \label{eq:example}
  E = mc^2 .
\end{equation}

As shown in Equation~\eqref{eq:example}, cross-references resolve after the
second compilation pass --- SciPen handles the reruns for you.

\section{Results}
\label{sec:results}

\begin{table}[ht]
  \centering
  \begin{tabular}{lrr}
    \toprule
    Condition & Accuracy & Time (s) \\
    \midrule
    Baseline  & 0.81     & 12.4     \\
    Ours      & 0.93     & 10.1     \\
    \bottomrule
  \end{tabular}
  \caption{Replace with your own results.}
  \label{tab:results}
\end{table}

\section{Conclusion}
\label{sec:conclusion}

Summarise the contribution and name the most promising next step.

\bibliographystyle{plain}
\bibliography{references}

\end{document}
